% GENERATED by md2tex.py from ../draft_v5_en/body/09_sim_to_real_gap.md -- do not edit by hand.
% Change the Markdown and run `python md2tex.py`; edits here are overwritten.

\chapter{Measuring the sim-to-real gap}
\label{ch:09}

\section{The gap as a subject of study}
\label{sec:9.1}

Adaptation A5 says that operational validation must include, as a required part and not an optional one, an explicit description of the gap between the environment where the system was trained and the one where it will run. This answers an imbalance in the literature pointed out in \S\ref{sec:2.7}: there are many techniques to \emph{reduce} the gap, but few frameworks to \emph{measure} it in a way a safety argument can use.

This chapter studies the gap in steps of increasing realism: the simulator where all the verdict evidence comes from, a bridge with more realistic physics and graphics, and the real platform. Each step answers a different question, and the chapter states clearly which questions have a measured answer today and which do not.

\section{First step: the high-fidelity bridge}
\label{sec:9.2}

An equivalent environment was built in Isaac Sim, which uses the PhysX physics engine and RTX rendering and is more realistic than Gazebo in both respects. The full chain was ported: import of the vehicle model, sensor publishing, in-process training, domain randomization, and the same cage with its parameter file.

The most useful result of this step is not a campaign but a diagnosis. It should be reported as such, because it is exactly the kind of information A5 is meant to collect. Porting the system uncovered a series of problems that were not about implementation but about how the layers depend on each other. The yaw authority of the model had to be recalibrated because the dynamic response is different. The cage's heading thresholds had to be recalibrated for the new renderer, because the vision estimator reads the same lane with a different bias. And the policy with speed control found a degenerate optimum (stopping) that the earlier environment had not allowed it to show.

The main lesson of this step comes from those problems: checkpoints do not transfer between simulators. A policy trained in one is not a valid policy in the other, even with the same architecture and the same reward. Each environment needs its own training. If this happens between two simulators, expectations about moving to the physical platform should be set with that in mind. It is a negative result, and a useful one.

\section{Second step: the physical platform}
\label{sec:9.3}

\subsection{Status}
\label{sec:9.3.1}

\begin{quote}
\textbf{What follows is the \emph{bring-up} of the physical step, not its results campaign}\sidenote{Bring-up: getting the hardware chain to run end to end and calibrating it. It comes before, and is not the same as, a campaign run under the scenario protocol.}, and this should be kept in mind before reading the numbers. The physical evidence in this chapter comes in two kinds, which carry different weight. The first kind is calibration measurements and structural findings: measurements against a reference, a regression over 5,665 cycles, two controlled pairs, a threshold comparison, and a true-position capture over the whole layout. These are results and do not depend on any campaign. The second kind is driving figures, which come from runs done outside the scenario protocol and are therefore preliminary: nominal tracking from a single run (\S\ref{sec:9.3.4}), and envelope blocking from three runs (\S\ref{sec:9.3.7}). A physical campaign would replace the second kind, not the first.
\end{quote}

The physical deployment chain is built, documented end to end and running on hardware: onboard camera, policy node, lane estimator, cage node, vehicle control and platform driver, all on the real vehicle on a closed lane circuit. This step has produced evidence of three kinds that should not be mixed: calibration measurements on the hardware, a negative result about the policy that simulation had validated, and a positive result about the policy retrained for transfer.

What it has not produced is a scored scenario, for three reasons that are independent of each other. The driving was done outside the library protocol, in monitoring mode instead of enforcement, and with a perception setting different from the one used to score the campaigns. Any one of these is enough to mean that nothing in this chapter is a verdict, and all three apply. So Chapter~\ref{ch:10} marks its physical column as \emph{not executed} (\S\ref{sec:10.4}d–\allowbreak{}e), and everything that follows here is later evidence. That column should not be confused with the physical column of the gap table in \S\ref{sec:9.4}, which does contain figures. The first one scores requirements under protocol and is empty by decision. The second one compares driving quantities outside the protocol and is preliminary by definition.

\subsection{The blocking check was done, and it failed}
\label{sec:9.3.2}

The item marked as blocking was the effective field of view of the onboard camera.\sidenote{The effective field of view is the horizontal angle the camera really sees. Converting pixels to millimetres depends on it, so a wrong value scales every lateral quantity.} It was a default configuration value that the simulation had copied, so no simulation result could reveal an error in it, and every lateral quantity the cage uses depends on it. It was measured before any driving metric, as the chapter had planned, and the value was wrong: the effective horizontal field of view of the real module is 77.89°, not the assumed 90°.

The effect was then measured against a tape measure, with the vehicle moved by hand over a range of $\pm$100 mm, on the unrectified image. The estimator read the lateral error as about 0.68–0.83 times the real value minus 10 mm. This under-reading did not go away with any filtering, and it would have triggered the boundary threshold, set at 160 mm, with the vehicle at a real 207–241 mm. The lane width, on the other hand, was read correctly (252.9 mm against a 250 mm ruler). A width is a \emph{difference} measured across the optical axis, while the lateral error is an \emph{absolute position} away from it, and the lens distortion term that was not modelled only compresses the second one. This diagnosis decided the fix, and the conclusion was clear: rectify the image, do not just change parameters.

\begin{quote}
\textbf{Partial retraction, which supports the chosen fix.} The numbers in the previous paragraph describe the \emph{unrectified path} and not the deployed system. When the measurement was repeated on the rectified image (nine points, vehicle on the ground, not touched), the scale is 1.058 on the left and 0.991 on the right, and the intercept is gone. The 160 mm threshold now triggers with the vehicle at a real 151–158 mm, which leaves about 100 mm of margin to the edge instead of 14 to 48. The under-reading was real, it came from the optical path and not from the estimator, and rectification removed it. The original number is kept because it is the evidence that justified the decision. No cage rule should be tuned from it.
\end{quote}

This is exactly the kind of finding adaptation A5 exists for, and it should be said clearly. An assumption that simulation could not falsify, because simulation had inherited it, turned out to be wrong on the very first measurement on hardware, and it sat in the path of a safety rule. The chosen fix is not to change the estimator's parameters but to rectify the real image to the standard camera model\sidenote{Rectifying removes the lens distortion, so that the image follows the same ideal pinhole camera model that the simulator uses.}, so that one projection works for both the simulation campaigns and the vehicle. Its effect was measured on hardware with the car stopped and untouched: cycles with invalid perception drop from 45\,\% to 5.5\,\%, and the lane boundary rule goes from firing 102 times without the car moving to not firing at all. The nine-point measurement above confirms this in the range that matters.

\subsection{The policy validated in simulation does not transfer}
\label{sec:9.3.3}

The first real drive with the reference policy, the two-dimensional one that the verdict campaign had evaluated, gave a clear negative result. The vehicle was moved by hand over a range of 332 mm of lateral error, with the chain running but not driving the motors, over 5,665 logged cycles. The steering command reacts to the position with the correct sign, but ten times too weakly. The change in steering across the whole 332 mm is 0.055, while there is a constant bias to the left of 0.1155, which is 2.1 times that whole change. Only 29 of the 5,665 samples (0.5\,\%) ask for a right turn. In closed loop the bias wins, and the vehicle leaves the lane to the left wherever it starts.

The cause of the bias is the training circuit itself. Driven clockwise, each lap has about 13 m of left curves and 2 m of right curves, a ratio of 6.5 to 1, and the training action log shows the mean steering flat at +0.112\ldots{}+0.120 over all 284,672 steps. It is not a perception or dynamics problem. The policy memorised the handedness of the track as a steering bias.\sidenote{Handedness is the balance between left and right curves. Driven clockwise, \texttt{complex\_b} has about 13 m of left curves and 2 m of right curves per lap.} This is a property of the training distribution, and no appearance randomization would have fixed it.

\subsection{Retraining for transfer, and the result}
\label{sec:9.3.4}

The gap was split into three separate terms, and only the first one is a property of the track. Handedness: observation and action are mirrored per episode\sidenote{Mirroring flips the image left to right and the sign of the steering together, so in about half of the episodes the policy effectively drives a right-handed circuit.}. Photometry: three quarters of the episodes use the lighting range measured in the hall. Camera geometry: mount pitch, height and, in one tenth of the episodes, the real measured lens. A new policy was trained on this basis for 2.5 million steps.

The deployed checkpoint is not the one with the highest reward. It was chosen on two criteria: how well it responds under deployment conditions, and how little it depends on the cage. The chosen candidate triggers the rate limiter in 3.0\,\% of nominal cycles, against 35.0\,\% for the reward peak. The second criterion was chosen on purpose. The transfer risk stated in Chapter~\ref{ch:08} was exactly the policy's dependence on the limiter, so a checkpoint that relies on it is the wrong one to put on hardware. This is the second time in this work that the reward curve picks badly.

The handedness term was fixed: the bias/excursion ratio drops to 0.07–1.10, against 12.9–19.2 for the reference policy as it was deployed.

\textbf{And it transfers: first measurement, one run.}\prelimflag{} The retrained policy drove 18.05 m of the real circuit in one continuous segment of 101.1 seconds, with no operator help. Its median lateral error was 18.7 mm (maximum 98.7 mm, against a boundary threshold of 160) and its maximum heading error was 18.91°, against a limit of 25. No safety rule fired during the run. The only rule that acted was the rate limiter, in 3.4\,\% of moving cycles, against 3.0\,\% measured in simulation.

That last pair of numbers is the most important observation of the bring-up, because it is the first measurement of a question Chapter~\ref{ch:08} left open. The transfer risk identified in advance was that the physical actuator does not apply the per-cycle change limit that the policy had come to rely on, so the pair \emph{(policy, limiter)} might not survive on hardware. In the first run it did, with the two numbers only four tenths of a percentage point apart. One run does not rule out the risk, and the campaign could still show it. But the difference is not borderline, and it goes in the opposite direction from the one feared. The framework asked for the risk to be stated before touching hardware and for the first measurement to test it, and that is what happened. Closing it needs the campaign.

\subsection{What stops the vehicle is not the policy}
\label{sec:9.3.5}

The run ended 2.11 m before completing the circuit, and the reason had nothing to do with driving. A single 400 ms pulse of invalid perception, which the estimator recovered from on its own, latched the emergency stop\sidenote{Latched: once \cagerule{05} fires it stays active until an explicit reset, even after the trigger has gone. This avoids switching on and off at the edge of the trigger.} while the vehicle was 27 mm from the lane centre, on the tightest curve of the layout. Because the rule's exit is asymmetric, it stayed latched for the remaining 396 s. The rule did exactly what its specification says: to exit, it needs an explicit reset. That decision came from the hazard analysis, to avoid switching back and forth at the edge of the trigger. The gap is not in the cage. It is between an artefact validated on simulated episodes, which end, and a vehicle that has to keep running. It was solved by leaving the rule as it is and placing the reset path outside the cage. This was done on purpose, so as not to change the artefact under verification with a change whose whole effect is on hardware and that simulation could not validate, because in simulation the latch hardly ever matters.

There are two more findings of the same kind, and neither involves the policy. When the stereo camera's visual odometry closes a loop\sidenote{Loop closure: when the camera recognises a place it has already seen, it corrects its accumulated drift in a single step. That step is what reaches the cage as a pose jump.}, it resets its pose and jumps. Since the state filter fuses that pose and not its velocity, the jump reaches the cage as a velocity spike. A displacement of 3.6 m in one frame and a velocity of 5.5 m/s were measured, against a contract speed of 0.22. This triggers the chain of speed rules and ends the run. Turning off loop closure removes the problem: jumps per frame go from 116 to 0 in 509 s. The second finding is that the cage's speed ceiling cannot fire on commanded motion in the deployed configuration, because its threshold (0.25 m/s) is higher than the operating speed (0.22 m/s). What Chapter~\ref{ch:08} recorded as a coverage gap in the campaign is, on hardware, a rule that does not protect the real case: the vehicle enters the tightest curve at contract speed with no speed protection from the cage. The audit of the caged runs adds the opposite problem, which is even more uncomfortable. The rule did fire, 58 and 40 cycles in the two runs that passed the audit, and 100\,\% of those were caused by velocity errors from the pose sensor (0.25–1.30 m/s reported, which the vehicle cannot reach under its own power). In one of the runs, these false activations blocked the reset path of the very rule they had triggered. A rule that does not act on the case that exists, but does act on one that does not, is a different finding from a coverage gap. The failure is not in the rule itself but in how its threshold, the chosen operating point and the quality of its only speed input fit together. It is documented and not fixed, because turning it back on would first need a reliable curvature signal (\S\ref{sec:9.3.6}).

\subsection[How accurate the estimator is depends on where the car is, not on how it moves]{How accurate the estimator is depends on \emph{where} the car is, not on how it moves}
\label{sec:9.3.6}

With the geometry fixed, the estimator still kept stopping the vehicle. The cause was found with a deliberately low-tech measurement: four stations marked with tape on the floor, the vehicle pushed by hand with a pointer on the chassis following the painted line, and only the camera running (no policy, no cage, no odometry). There were four laps and four tests with the vehicle stopped, about 11,500 frames in total. The acceptance criterion was set before measuring and is based on a geometric fact that does not depend on odometry: on a closed circuit, the curvature integrated over one lap equals 2$\pi$\sidenote{Over one closed lap the heading turns through one full circle, whatever the shape of the circuit. The check needs no reference position, and the measured values are ratios to that expected total.}.

The criterion fails by a factor of two. The integration gives 1.97 / 2.01 / 2.25 / 2.37, against an acceptable range of 0.75–1.35. This is the first measurement of the curvature channel that does not depend on the vehicle or its state estimation. But the result that changes the order of the work is a different one: the estimator does not get worse because of driving, it gets worse depending on where it is. At the start of the straight, it pairs 96.7\,\% of the frames and is off by 7.2 mm. On the sections in between, it pairs between 37.9\,\% and 60.2\,\%. With the vehicle stopped at two of those points and the true lateral error at zero, one test pairs 0.0\,\% of its 273 frames, and another pairs 100\,\% with a bias of $-$39.7 mm and a standard deviation of 3.1 mm. In other words, it is confidently wrong, so no check based on spread can catch it.

This has two consequences for the method, and both change what was written earlier. The first is that the start of the straight, where the field of view calibration, the nine-point measurement and all the pre-deployment checks were done, turns out to be the best point of the circuit for the estimator. The content of those measurements still holds, but they can no longer be taken as valid for the whole layout. The second is about the mechanism. The problem is not choosing the wrong pair among candidate lines. It is in how the candidates are generated: edges of the same stripe, nearby markings, or a single visible line. The pair that survives has a plausible width, but its midpoint is shifted. A check for continuity over time on the chosen pair would not fix this, because the wrong pair is stable to within 3.1 mm.

\subsection{The three cage blockages, measured}
\label{sec:9.3.7}

The claim in \S\ref{sec:9.3.5}, that what stops the vehicle is not the policy, was tested in three runs in a row on the real circuit, all in monitoring mode, changing only the settings of the reset path. None of them kept the vehicle moving, and all three failed because of measurement, not driving.

The first run shows that the requirement of one second of healthy perception cannot be met while moving. Nine resets were issued, and the log contains 623 that were denied. 48\,\% of those were denied because the wait never completed, since the estimator's health signal flickers faster than once per second. That threshold came from a hazard analysis argument against switching at the edge of the trigger, and it had never been tested with the vehicle moving. The second run tests the documented way out (removing the lane boundary rule from the blocking set) and rules it out. The budget of thirty resets is used up in about one minute, each one latching again immediately, with the vehicle stopped and the estimator reading a lateral error of $-$296 mm. A denied reset just becomes a wasted reset. It is not a way out. The third run shortens the wait and finds the next thing that blocks: over the 1,066 cycles in which it acts, the heading rule fires with the vehicle 20 mm from the lane centre and a heading standard deviation of 19.1°, against a threshold of 25°. The same configuration, measured at the start of the straight, gave 5.3°.

That last comparison is the important one. It is the fourth independent confirmation of \S\ref{sec:9.3.6}, and the first in the heading channel with the vehicle driving: heading gets worse away from the good spot, in the same way as offset and curvature. The three blockages do not show a policy that cannot drive. They show the stopping rule firing on a measurement that fails exactly where \S\ref{sec:9.3.6} said it fails. That is why the remaining work puts the estimator before the policy.

\subsection{Experimental design still to do}
\label{sec:9.3.8}

The physical subset of the library (nominal, and nominal with curvature and heading limit, in both operating modes) has still not been run under protocol, and four conditions stand in the way. Three were identified before driving and are about instrumentation and contract: a reset path for the emergency stop, per-run provenance logging (already implemented), and a decision on which perception contract to use, because the heading setting that lets the vehicle drive is not the one the campaigns were scored with. They are listed with their criteria in \S\ref{sec:12.4} (T2), which is where the work hands them over.

What should be stated here is which of the four matters most, because it is a result of this step and not an inherited task. The fourth one was added by the hardware: widen the lane estimator before restricting the policy. The estimator's accuracy depends on the point of the circuit (\S\ref{sec:9.3.6}), and its failure is what actually blocks the vehicle (\S\ref{sec:9.3.7}). Until this is fixed, a scored run would measure the envelope on an input that is known to be faulty in areas the vehicle normally drives through. Only one parameter of the operational domain is meant to be closed at this step, and only here: the lateral acceleration envelope, which cannot be measured in simulation at all. It is still open.

\section[The gap table (preliminary physical column)]{The gap table \emph{(preliminary physical column)}}
\label{sec:9.4}

The table now has a physical column\prelimflag{}, and the first thing to notice is not a row but a column heading: the deployed policy is not the one that produced the simulation verdict. That change is itself a result of this step (\S\ref{sec:9.3.3}). It means the physical column and the simulation columns are not a like-for-like comparison. They describe different policies, in different modes, under different perception contracts. They are shown together because the comparison still tells something useful, and they are separated in their headings because treating them as a clean difference would be exactly the kind of unsupported claim the framework is meant to prevent.

% layout: table*, fixed, \small, 164 of 165 mm
\begin{table*}[htbp]
\caption[Gap table. The two simulation columns describe the verdict policy.]{Gap table. The two simulation columns describe the verdict policy. The physical column is preliminary (the policy retrained for transfer, in monitoring, over one run, with no scored scenario) and the physical campaign will replace it.}\label{tab:9.1}
\small\setlength{\tabcolsep}{4pt}
\begin{tabular}{@{}>{\RaggedRight\arraybackslash}p{24.9mm}>{\RaggedRight\arraybackslash}p{27.9mm}>{\RaggedRight\arraybackslash}p{39.1mm}>{\RaggedRight\arraybackslash}p{63.2mm}@{}}
  \toprule
  \tabhead{Metric (nominal scenario)} & \tabhead{Sim 1-D \emph{(historical reference)}} & \tabhead{Sim 2-D \emph{(the one that produced the verdict)}} & \tabhead{\textbf{Physical} \emph{(PRELIMINARY: retrained policy, monitoring, \textbf{one run}, unscored)}} \\
  \midrule
  Median lateral error & 10.9 mm & 8.6 mm (max. 27.3 mm) & 18.7 mm (max. 98.7 mm) \\
  \addlinespace[2pt]
  Continuous run & 4.88 laps & 5.32 laps & 18.05 m in one segment ($\approx$ 0.94 of the perimeter) \\
  \addlinespace[2pt]
  Maximum lateral excursion & < limit over the whole domain & 0 edge contacts inside the ODD & 0 contacts; max. 98.7 mm against a threshold of 160 \\
  \addlinespace[2pt]
  Emergency stops & 0 & 0 & 1, from perception, with the vehicle in the lane \\
  \addlinespace[2pt]
  Intervention rate & 43.5\,\% (limiter only) & 3.0\,\% at the deployed checkpoint & 3.4\,\% (limiter only) \\
  \addlinespace[2pt]
  Operating speed & 0.200 m/s fixed & $\approx$ 0.216 m/s under a 0.22 ceiling & $\le$ 0.213 m/s under the same ceiling \\
  \addlinespace[2pt]
  Control loop rate & 10 Hz nominal & 10 Hz nominal & 8.68 Hz measured; 9.6 Hz after removing the lidar from layer 2 \\
  \bottomrule
\end{tabular}
\end{table*}

Three observations. First, the row that held the most risk no longer does. Chapter~\ref{ch:08} named the limiter's intervention rate as the first gap to watch, because what transfers is not the policy alone but the pair of policy and limiter, and the physical actuator does not apply that limit. With the checkpoint chosen for independence from the cage, the two numbers differ by four tenths of a percentage point. The risk was correctly identified, and the way it was handled (choosing by independence, not by reward) worked.

\textbf{The lateral error doubles, which was expected.} It is 18.7 mm against 8.6 mm in simulation, with a loop running at 8.68 Hz instead of the 10 Hz the policy was trained for, and with an estimator whose accuracy depends on the point of the circuit (\S\ref{sec:9.3.6}). It stays well within the cage threshold, and the number to watch is the 98.7 mm maximum, not the median. It cannot be claimed that this maximum hides a larger real value because the estimator under-reads: on the rectified path the scale is practically one (\S\ref{sec:9.3.2}). The bias that does remain is the one that depends on place. It can go in either direction, and without true-position measurements it cannot be applied row by row.

\textbf{The loop rate is still open, but it is no longer the main gap.} It is the only parameter of the deployment contract that the hardware does not reach. Its cause has been found: not the camera, which delivers 9.84 Hz, but the network's inference time. It has been improved from 7.3 to 9.6 Hz, but not closed.

\section{Summary}
\label{sec:9.5}

The chapter leaves the gap clearly described, and in a different state from before. The first step is measured, and its result is still a useful negative diagnosis: checkpoints do not transfer between simulators, and perception, calibration and dynamics depend on each other more than the modular architecture suggests. The second step ends with the bring-up done and the results campaign not done, and it is better to say that directly than to present it as a measurement in progress. Even so, the bring-up produced five things that simulation could not, and four of them do not depend on the campaign being run.

It produced a falsified assumption: the field of view that simulation had inherited from a default setting was wrong, it sat in the path of a safety rule, and no amount of extra simulation would have shown it. It produced a negative result about the validated policy: it does not transfer, and the cause, the handedness of the training circuit memorised as a steering bias, is a property of the training distribution and not of the hardware. It produced a positive result about the replacement policy, and with it an answer to a transfer risk stated in advance. It located a failure mode that simulation cannot show at all: the lane estimator's accuracy depends on the \emph{place} on the circuit and not on what the vehicle is doing. This was measured with true-position methods and later confirmed independently in the heading channel while driving. And it produced a kind of finding that only appears with a real vehicle in front of you: rules that are correctly specified but have no defined behaviour in operation, thresholds that cannot fire in the deployed configuration, and sensors whose failures go straight into the cage's only speed input.

There is also a result about the method, which only becomes visible when all the gap terms are put into one table. This work can state it because it wrote its list of gaps \emph{before} touching hardware and did not edit it afterwards. That early list got four terms right (appearance, provisional thresholds, computation rate and yaw gain), and they are all of the same type: terms that a simulator knows how to represent. Nothing that actually stopped the vehicle was on that list: the handedness of the training circuit memorised as a bias, an intrinsics error that the simulator had inherited and so could not reveal, an estimator whose accuracy depends on place and not on action, a correctly specified latch with no defined operational behaviour, and a sensor whose failures go straight into the cage's only speed input. Three of these five are in the measurement chain (the camera intrinsics, the lane estimator and the speed sensor), one is a property of the training distribution, and one is about how a rule behaves in operation. None of the five is the control policy. This imbalance is the chapter's methodological finding: the gap terms a team working in simulation can list in advance are exactly the ones simulation knows how to model (appearance, timing, gains), and the ones that stopped this vehicle are the ones simulation had no representation of at all, either because it inherited the same assumption, because the quantity does not exist there, or because its plant is better than the real one.

Adaptation A5 is met for the first half of its requirement, which is to make the gap explicit and organise how it is measured. The second half has started but is not finished: there are measurements on hardware, but from bring-up and not from a campaign, so there is no per-scenario gap table and no physical verdict. That, and nothing else, is why the verdict column in Chapter~\ref{ch:10} is marked not executed. Chapter~\ref{ch:10} gives the final verdict with this limitation built into its statement, and with one clarification it could not make before: the limitation is no longer \enquote{there is no physical evidence} but \enquote{the physical evidence comes from bring-up and is not scored}.

There is one more point that none of the above covers, and Chapter~\ref{ch:10} includes it in its statement: the envelope has never acted on the real vehicle. All physical runs were done in monitoring mode, with the safe action equal to the raw action in every logged cycle. The central artefact of this work is verified in simulation and only \emph{observed} on hardware, and the campaign that would change that is the first item of future work.
